\section{The truth order and truth-order closure}

\begin{proposition}[The truth order as a relation]\label{thm:tle}\lean{Obsmetrologia.Reading.TLE}\leanok\uses{thm:four-bilattice}
The truth order is the relation \texttt{TLE}, the \texttt{Prop}-valued form
of the Boolean comparison \texttt{leT}. It is reflexive, transitive, and
antisymmetric by finite charge analysis.
\end{proposition}

\begin{theorem}[$\vee_t$ and $\wedge_t$ are the truth lub and glb]\label{thm:truth-lattice}\lean{Obsmetrologia.Reading.joinT_le}\leanok\uses{thm:tle}
Truth join is the least upper bound and truth meet is the greatest lower
bound for \texttt{TLE}. With $f$ least and $t$ greatest, this is the bounded
truth lattice of FOUR.
\end{theorem}

\begin{theorem}[Negation is the order duality]\label{thm:negate-tle}\lean{Obsmetrologia.Reading.negate_tle}\leanok\uses{thm:truth-lattice, thm:negate-duality}
As a relation, negation reverses the truth order and preserves the knowledge
order. This is the order-dual reading of the same slot swap.
\end{theorem}

\begin{proposition}[Truth-order closure operator]\label{thm:tclosure}\lean{Obsmetrologia.Reading.TClosure}\leanok\uses{thm:tle}
A truth-order closure is an operator on charges that is inflationary,
monotone, and idempotent with respect to \texttt{TLE}.
\end{proposition}

\begin{proposition}[The canonical truth closure and truth-settling]\label{thm:settleT}\lean{Obsmetrologia.Reading.joinTClosure}\leanok\uses{thm:tclosure, thm:truth-lattice}
The canonical truth closure raises a charge toward a chosen truth bound by
truth join. Truth-settling is the fixed point reached by this closure.
\end{proposition}

\section{The light cone}

\begin{theorem}[FOUR is the $1{+}1$ causal diamond in null coordinates]\label{thm:light-cone}\lean{Obsmetrologia.Reading.time, Obsmetrologia.Reading.space, Obsmetrologia.Reading.kle_iff_null, Obsmetrologia.Reading.tle_iff_null, Obsmetrologia.Reading.kle_iff_lightcone, Obsmetrologia.Reading.kle_time_mono, Obsmetrologia.Reading.kBot_past_tip, Obsmetrologia.Reading.kTop_future_tip, Obsmetrologia.Reading.tTop_tBot_spacelike, Obsmetrologia.Reading.coords_injective}\leanok\uses{thm:four-bilattice, thm:kBot-least, thm:tle, thm:negate-tle}
Read the two evidence slots $(\mathrm{pos},\mathrm{neg})$ as the null coordinates
$(u,v)$ of a two-dimensional Minkowski diamond. The knowledge order is the product
order on $(u,v)$, which is the causal-future relation in null coordinates
(\texttt{kle\_iff\_null}); equivalently $a\le_k b$ iff
$|\mathrm{space}\,b-\mathrm{space}\,a|\le \mathrm{time}\,b-\mathrm{time}\,a$
(\texttt{kle\_iff\_lightcone}), where $\mathrm{time}=u+v\in\{0,1,2\}$ is the
knowledge content and $\mathrm{space}=u-v$. The truth order raises $u$ and lowers
$v$: the spacelike direction (\texttt{tle\_iff\_null}). $\bot$ is the past tip and
$\top$ the future tip; $t$ and $f$ are at equal time, $\le_k$-incomparable and
$\le_t$-comparable, i.e.\ spacelike separated. The monograph's ``one diamond turned
a quarter turn'' is the change of coordinates from $(u,v)$ to
$(\mathrm{time},\mathrm{space})$. Every declaration is axiom-free: the statements
are finite checks over the four charges, which is exactly what a $1{+}1$ diamond
with one bit per null ray is.
\end{theorem}

\begin{theorem}[Negation is parity, conflation is time reversal: $K_4=\{1,P,T,PT\}$]\label{thm:light-cone-k4}\lean{Obsmetrologia.Reading.time_negate, Obsmetrologia.Reading.space_negate, Obsmetrologia.Reading.time_conflate, Obsmetrologia.Reading.space_conflate, Obsmetrologia.Reading.act_centered, Obsmetrologia.Reading.act_space, Obsmetrologia.Reading.signs_injective, Obsmetrologia.Reading.timeSign_mul, Obsmetrologia.Reading.spaceSign_mul, Obsmetrologia.Reading.k4_is_PT}\leanok\uses{thm:light-cone, thm:klein-four-symmetry, thm:k4-group}
Negation fixes time and reverses space; conflation reflects time through the middle
of the diamond and fixes space. On the centred coordinates $(\mathrm{time}-1,
\mathrm{space})$ the Klein four group acts as the sign-change group $\{1,P,T,PT\}$,
faithfully (\texttt{signs\_injective}) and multiplicatively
(\texttt{timeSign\_mul}, \texttt{spaceSign\_mul}). The external identification,
stated and not formalized: $\{1,P,T,PT\}$ is the component group of the Lorentz
group $O(1,1)$, and equally of $O(3,1)$; the gauge of FOUR is the discrete part of
any Lorentzian symmetry group, compatible with $3{+}1$ without selecting the $3$.
\end{theorem}

\begin{theorem}[Chu duality is the transpose composed with parity]\label{thm:light-cone-chu-dual}\lean{Obsmetrologia.Reading.chu_dual_time, Obsmetrologia.Reading.chu_dual_space}\leanok\uses{thm:light-cone-k4, thm:predication-chu}
The dual of a FOUR-valued Chu space transposes the matrix and negates each entry;
in light-cone coordinates it therefore fixes every entry's time and reverses its
space. The space-for-time transpose of the monograph's spatiotemporal Chu space is
this duality read on the diamond.
\end{theorem}
